\chapter{Applications of JSONPath in Data Analytics and Natural Language Processing}
\label{chap:applications-jsonpath}

\section{Introduction}

JSONPath is widely used in data analytics pipelines for filtering and transforming structured data, and in NLP workflows for extracting annotations or metadata from JSON-based formats \cite{jsonpath_data_analytics, restfulapi_jsonpath}.

\section{Use Cases in Data Analytics}

\begin{itemize}
	\item \textbf{ETL and Data Integration}: JSONPath expressions are very impotent in ETL workflows, allow accurate filtering, transformation, and extraction of JSON data during ingestion and preparation stages \cite{jsonpath_data_analytics}.
	\item \textbf{Streaming Analytics}: Platforms like Amazon Kinesis Data Analytics uses JSONPath to map and extract nested JSON fields (e.g., customer order data), enabling real-time processing and analysis \cite{kinesis_jsonpath}.
\end{itemize}

\section{Applications in Natural Language Processing (NLP)}

\begin{itemize}
	\item \textbf{Standardized Interchange Formats}: The JSON-NLP schema gives a unique JSON for NLP annotations (like tokens, entities, syntactic structures), making JSONPath a practical tool for extracting annotation layers in a well structured and automated way \cite{json_nlp_schema}.
	\item \textbf{Preprocessing and Annotation Extraction}: In NLP pipelines, JSONPath allows developers to select and retrieve annotation fields—like named entities or syntactic metadata—from deeply nested JSON outputs, improve flexibility and simplify code.
\end{itemize}

\section{Summary Table}

\begin{tabular}{p{5cm} p{9cm}}
	\hline
	\textbf{Scenario} & \textbf{Role of JSONPath} \\
	\hline
	Data integration / ETL & Filters and transforms JSON records during extraction and ingestion (e.g., sales totals, nested records) \cite{jsonpath_data_analytics} \\
	Streaming analytics & Maps streaming JSON sources to SQL columns in real-time analytics pipelines (e.g., through Amazon Kinesis) \cite{kinesis_jsonpath} \\
	NLP annotation extraction & Retrieves structured linguistic annotations from JSON-NLP formatted outputs effectively \cite{json_nlp_schema} \\
	\hline
\end{tabular}

\section{Conclusion}

JSONPath fills the gap between structured JSON data and analytical/NLP workflows by providing a simple way to query and extract relevant information. Declarative syntax in JSONPath makes it possible for data engineers and NLP researchers to integrate JSON pipelines without clunky imperative traversal code.